\chapter{Peritoneoscopy}
\index{Peritoneoscopy}

\medskip
\noindent\textit{This chapter is for general education. The treating surgical team should explain the reason for the procedure, alternatives, and individual risks.}

\section*{What it is}
Peritoneoscopy is a minimally invasive procedure that uses a camera through a small abdominal incision to inspect the peritoneal cavity and abdominal organs. It is often performed as part of laparoscopy and may allow biopsy or treatment during the same procedure.

\section*{Why it may be done}
It can help investigate unexplained abdominal or pelvic symptoms, assess disease such as endometriosis or cancer, inspect the liver or peritoneum, obtain tissue, or guide selected operations. Imaging or less invasive tests may be appropriate alternatives depending on the question.

\section*{Before and during the procedure}
The team will review medicines, allergies, pregnancy possibility, previous operations, and anaesthetic risks. You may need fasting and instructions about medicines. The procedure is usually done with anaesthesia; carbon dioxide may be used to create working space. A specimen may be sent to a laboratory.

\section*{Recovery and risks}
Short-term shoulder discomfort, bloating, nausea, tiredness, and pain around the incisions can occur. Risks include bleeding, infection, injury to bowel, bladder, blood vessels, or other organs, blood clots, anaesthetic complications, and the need to convert to open surgery. Recovery depends on whether treatment was also performed.

\section*{When to Seek Medical Care}
After the procedure, seek urgent care for severe or increasing abdominal pain, persistent vomiting, fever, fainting, breathing difficulty, chest pain, heavy bleeding, a swollen abdomen, or redness and discharge from an incision.

\section*{Sources and further reading}
\begin{itemize}
\item MedlinePlus. \textit{Laparoscopy}. \url{https://medlineplus.gov/ency/article/003918.htm}
\item National Institute of Diabetes and Digestive and Kidney Diseases. \textit{Laparoscopic surgery}. \url{https://www.niddk.nih.gov/health-information/diagnostic-tests/laparoscopy}
\end{itemize}
